\documentclass[11pt,letterpaper]{article}
\usepackage[T1]{fontenc}
\usepackage{lmodern,microtype}
\usepackage[margin=.75in,headheight=15pt]{geometry}
\usepackage{amsmath,amssymb,mathtools}
\usepackage{xcolor,enumitem,fancyhdr,hyperref}
\definecolor{ink}{HTML}{222834}
\definecolor{indigo}{HTML}{25355D}
\definecolor{muted}{HTML}{505A6B}
\definecolor{keyink}{HTML}{205D52}
\color{ink}
\hypersetup{colorlinks=true,linkcolor=indigo,urlcolor=indigo,
  pdftitle={Number Theory Makeup Midterm},pdfauthor={Igor Rivin}}
\setlength{\parindent}{0pt}
\setlength{\parskip}{4pt}
\setlength{\abovedisplayskip}{6pt}
\setlength{\belowdisplayskip}{6pt}
\clubpenalty=10000
\widowpenalty=10000
\setlist[enumerate]{label=(\alph*),leftmargin=1.7em,
  itemsep=4pt,topsep=5pt,parsep=2pt}
\newcommand{\Prob}{\mathbb P}
\newcommand{\E}{\mathbb E}
\DeclareMathOperator{\Var}{Var}
\newcommand{\edition}{\ifanswers Instructor key\else Student paper\fi}
\newcommand{\partpts}[1]{\unskip\nobreak\hspace{.3em}%
  \mbox{\small\color{muted}(#1 points)}}
\newcommand{\credit}[1]{\par{\footnotesize\color{muted}\textit{Credit:} #1}}
\newcommand{\problem}[3]{\par\addvspace{9pt}%
  {\large\bfseries\color{indigo}#1.\ #2}\hfill
  {\small\color{muted}\textbf{#3 points}}\par\nobreak}
\NewDocumentEnvironment{solution}{O{1in}+b}{%
  \ifanswers
    \par\begingroup\small\color{keyink}\textbf{Solution.} #2\par\endgroup
  \else
    \par\nobreak\vspace*{#1}
  \fi
}{}
\pagestyle{fancy}
\fancyhf{}
\fancyhead[L]{\small\color{muted}Number Theory \textbar\ Makeup Midterm}
\fancyhead[R]{\small\color{muted}Patterns, Proofs, and Primes}
\fancyfoot[L]{\footnotesize\color{muted}\edition}
\fancyfoot[R]{\footnotesize\color{muted}Page \thepage\ of \pageref*{exam-last-page}}
\renewcommand{\headrulewidth}{.3pt}
\begin{document}
{\LARGE\bfseries\color{indigo}Number Theory Makeup Midterm}\par
{\small\color{muted}50 minutes \quad $\bullet$ \quad 100 points
  \quad $\bullet$ \quad Six problems}\par
\ifanswers
\textbf{Grading note.} Award method credit after earlier errors. Suggested
working times for Problems 1--6 are 7, 7, 7, 7, 6, and 10 minutes,
leaving 6 minutes to check. Coverage is based on the midterm review bank.
\else
\smallskip
\textbf{Name:}\ \rule{.49\linewidth}{.4pt}\hfill
\textbf{Date:}\ \rule{.20\linewidth}{.4pt}\par
\fi
\textbf{Instructions.} Answer all six problems and show your reasoning.
No calculator is needed. You may quote results proved in class unless a
proof is requested. Use the spaces provided; clearly label any work
continued on the back. Reserve time to check your work.

\problem{1}{Euclid and a linear congruence}{15}
\begin{enumerate}
\item Use the Euclidean algorithm to find $\gcd(195,84)$. \partpts{5}
\begin{solution}[.7in]
$195=2\cdot84+27$, $84=3\cdot27+3$, $27=9\cdot3$.
The last nonzero remainder is $3$, so the gcd is $3$.
\credit{Correct divisions: 4; identifying the gcd: 1.}
\end{solution}
\item Express this gcd as an integer linear combination of $195$ and $84$.
\partpts{5}
\begin{solution}[.65in]
$3=84-3\cdot27=84-3(195-2\cdot84)=7\cdot84-3\cdot195$.
\credit{Back substitution: 3; correct coefficients: 2. Accept any valid pair.}
\end{solution}
\item Find all solutions of $84x\equiv6\pmod{195}$. Give the distinct
residues modulo $195$ and explain why your list is complete. \partpts{5}
\begin{solution}[.8in]
Doubling (b) gives $x=14$ as a solution. Dividing the congruence by $3$
gives $28x\equiv2\pmod{65}$. Since $28$ is coprime to $65$, this has one
class modulo $65$, namely $14$. The three classes modulo $195$ are
$14,79,144$.
\credit{A solution: 2; all three classes: 2; completeness: 1.}
\end{solution}
\end{enumerate}

\problem{2}{Continued fractions and an inverse}{15}
Let $D(a)=\begin{psmallmatrix}a&1\\1&0\end{psmallmatrix}$.
\begin{enumerate}
\item Find the simple continued fraction of $139/60$ and all its
convergents. Use a final digit at least $2$. \partpts{6}
\begin{solution}[.85in]
Euclid gives remainders $19,3,1,0$ and quotients $2,3,6,3$.
Thus $139/60=[2;3,6,3]$, with convergents
$2/1,7/3,44/19,139/60$.
\credit{Digits: 2; convergents: 1 each.}
\end{solution}
\item Multiply the digit matrices for your expansion. Use the determinant
of their product to find the inverse of $60$ modulo $139$. \partpts{9}
\begin{solution}[1.05in]
\[
D(2)D(3)D(6)D(3)=\begin{pmatrix}139&44\\60&19\end{pmatrix}.
\]
The determinant is $(-1)^4=1$, hence $139\cdot19-60\cdot44=1$.
Therefore $60^{-1}\equiv-44\equiv95\pmod{139}$.
\credit{Matrix: 5; determinant/B\'ezout identity: 2; inverse: 2.}
\end{solution}
\end{enumerate}

\newpage
\problem{3}{Reassembling residues and reducing powers}{15}
\begin{enumerate}
\item Solve $x\equiv2\pmod5$ and $x\equiv4\pmod7$.
Justify uniqueness modulo $35$. \partpts{7}
\begin{solution}[1.25in]
Write $x=2+5k$. Modulo $7$, $5k\equiv2$, so $k\equiv6$.
Thus $x\equiv32\pmod{35}$. A difference of solutions is divisible by
both $5$ and $7$, hence by $35$, and adding $35$ preserves both residues.
\credit{Substitution and solution: 4; coprimality and uniqueness: 3.}
\end{solution}
\item Find the multiplicative order of $3$ modulo $7$ and use it to
compute $3^{100}\pmod7$. Show why your order is the least possible.
\partpts{8}
\begin{solution}[1.25in]
The successive powers $3^1,\ldots,3^6$ modulo $7$ are
$3,2,6,4,5,1$. The first occurrence of $1$ is at exponent $6$,
so the order is $6$. Since $100\equiv4\pmod6$,
$3^{100}\equiv3^4\equiv4\pmod7$.
\credit{Order with minimality justification: 5; exponent reduction and result: 3.}
\end{solution}
\end{enumerate}

\problem{4}{Prime factors and divisor sums}{15}
Here $\tau(n)$ counts positive divisors, $\varphi(n)$ counts units modulo
$n$, and $\mu$ is the M\"obius function.
\begin{enumerate}
\item Factor $180$ and compute $\tau(180)$, $\varphi(180)$,
and $\mu(180)$. Show the prime-factor calculations. \partpts{6}
\begin{solution}[1in]
$180=2^2 3^2 5$. Hence
\[
\tau(180)=3\cdot3\cdot2=18,\qquad
\varphi(180)=180\left(1-\frac12\right)\left(1-\frac13\right)
\left(1-\frac15\right)=48.
\]
Since a prime square divides $180$, $\mu(180)=0$.
\credit{Factorization: 1; divisor count: 2; totient: 2; M\"obius value: 1.}
\end{solution}
\item Suppose $f$ satisfies $\sum_{d\mid n}f(d)=n$ for every positive
integer $n$. Use M\"obius inversion to give a formula for $f(n)$ and
compute $f(12)$. \partpts{9}
\begin{solution}[1.2in]
Inversion gives
\[
f(n)=\sum_{d\mid n}\mu(d)\frac nd.
\]
At $n=12$, only $d=1,2,3,6$ contribute, so
$f(12)=12-6-4+2=4$.
\credit{Correct inversion formula: 5; nonzero terms: 3; value: 1.}
\end{solution}
\end{enumerate}

\newpage
\problem{5}{Valuations and a $5$-adic limit}{15}
For a prime $p$, let $v_p(r)$ denote the exponent of $p$ in a nonzero
rational $r$, with $v_p(a/b)=v_p(a)-v_p(b)$. Define
$|r|_5=5^{-v_5(r)}$ for $r\ne0$, $|0|_5=0$, and $d_5(x,y)=|x-y|_5$.
\begin{enumerate}
\item State Legendre's formula for $v_p(n!)$ and use it to compute
$v_3(30!)$. \partpts{5}
\begin{solution}[1.45in]
\[
v_p(n!)=\sum_{j\ge1}\left\lfloor\frac{n}{p^j}\right\rfloor,
\qquad v_3(30!)=10+3+1=14.
\]
Terms vanish once $p^j>n$.
\credit{Formula: 2; summands: 2; result: 1.}
\end{solution}
\item Compute $|125/14|_5$ and $d_5(1/3,16/3)$. \partpts{5}
\begin{solution}[1.45in]
Since $v_5(125/14)=3$, $|125/14|_5=1/125$.
Also $1/3-16/3=-5$, of valuation $1$, so the distance is $1/5$.
\credit{First valuation and norm: 2; difference, valuation, and distance: 3.}
\end{solution}
\item Let $S_n=\sum_{j=0}^n5^j$. Find its $5$-adic limit in
$\mathbb Q$ and justify your answer by computing the distance to the limit.
\partpts{5}
\begin{solution}[2in]
The geometric sum is $S_n=(5^{n+1}-1)/4$. Thus the limit is $-1/4$,
since
\[
d_5(S_n,-1/4)=\left|\frac{5^{n+1}}4\right|_5
=5^{-(n+1)}\longrightarrow0.
\]
\credit{Geometric sum and proposed limit: 2; exact distance: 2; convergence: 1.}
\end{solution}
\end{enumerate}

\newpage
\problem{6}{A field of formal pairs}{25}
Use formal pairs $a+bs$ with $a,b\in\mathbb F_7$ and the relation $s^2=3$.
Addition is coefficientwise; multiplication is distributive with $s^2$
replaced by $3$. You may use the resulting commutative ring laws.
All scalar arithmetic in this problem is modulo $7$.
\begin{enumerate}
\item List the squares in $\mathbb F_7$. Show that $a^2-3b^2\ne0$
whenever $(a,b)\ne(0,0)$, and deduce that every nonzero pair has inverse
$(a-bs)/(a^2-3b^2)$. \partpts{11}
\begin{solution}[1.65in]
The squares are $0,1,2,4$, so $3$ is nonsquare. If $b=0$, then $a\ne0$
and $a^2\ne0$. If $b\ne0$, vanishing of $a^2-3b^2$ would give
$(a/b)^2=3$, a contradiction. Since
$(a+bs)(a-bs)=a^2-3b^2$, division by this nonzero scalar supplies an inverse.
Thus the ring is a field.
\credit{Squares/nonsquare: 4; denominator proof (both cases): 5; inverse: 2.}
\end{solution}
\item Compute $(2+s)(3+2s)$ and $(1+s)^{-1}$. Check the inverse by
multiplication. \partpts{8}
\begin{solution}[1.5in]
\[
(2+s)(3+2s)=6+7s+2s^2=12=5.
\]
Also $(1+s)^{-1}=(1-s)/(1-3)=3(1-s)=3+4s$.
Checking: $(1+s)(3+4s)=3+7s+4s^2=15=1$.
\credit{Product: 4; inverse: 2; multiplication check: 2.}
\end{solution}
\item Give a basis over $\mathbb F_7$, the cardinality, and the
characteristic. Explain why this field differs from $\mathbb Z/49\mathbb Z$.
\partpts{6}
\begin{solution}[1.3in]
The basis is $1,s$: formal pairs have unique coefficients in
$\mathbb F_7$. There are $7^2=49$ elements and the characteristic is $7$.
In $\mathbb Z/49\mathbb Z$, the characteristic is $49$ and the nonzero
class $7$ has square zero, so that ring is not a field.
\credit{Basis and justification: 2; cardinality: 1; characteristic: 1;
valid distinction from the residue ring: 2.}
\end{solution}
\end{enumerate}
\vfill
\begin{center}\small\color{muted}End of paper. Check that you have answered all six problems.\end{center}
\label{exam-last-page}
\end{document}
